\chapter{Metodología}

En este capítulo se va a exponer cómo se ha desarrollado el proyecto en cuanto a organización y planificación de cada
una de las partes que lo han compuesto.

\section{Metodología de desarrollo}

Para la organización del proyecto se ha utilizado la metodología ágil \hyperlink{scrum}{\textbf{Scrum}} \cite{scrum}. 
Este tipo de metodología es muy útil debido a que permite ir realizando entregas del proyecto de forma iterativa y 
permite adaptarse a los cambios de forma rápida y eficiente.

A continuación se va a proceder a explicar en distintos apartados las piezas clave que componen el marco Scrum.

\subsection{Miembros del equipo Scrum}

Un equipo Scrum está formado por varios miembros, desempeñando cada uno un rol junto con unas responsabilidades.
Aunque normalmente la cantidad de miembros es reducida, pueden darse situaciones donde este número crezca considerablemente
porque así lo requiera el proyecto pero por lo general, los equipos de Scrum deben componerse de tres roles: el propietario
del producto (o \textit{Product Owner}), el \textit{Scrum Master} y el equipo de desarrollo.

En el proyecto, estos roles han sido asignados de la siguiente manera:

\begin{itemize}
    \item \textbf{Propieatrio del producto}
    \newline
    Este rol ha sido asignado tanto al tutor del TFG, el Dr. Juan Luis Jiménez 
    Laredo, así como al autor del presente proyecto. Se ha decidido así puesto que, entre ambos se ha ido dando forma
    a la propuesta: qué necesidades debía de cubrir el sistema, los requisitos funcionales  y las historias de usuario. 
    Asimismo, se encargaron de la priorización en cada iteración y de determinar cuando se consideraba cerrada una iteración,
    comprobando que el desarrollo fuese funcionando y incremental.

    \item \textbf{Scrum Master}
    \newline
    Este rol ha sido asignado al autor del proyecto, quien se ha encargado de asegurar el 
    cumplimiento de las prácticas Scrum mediante la orientación, eliminación de impedimentos, optimización del rendimiento 
    y entregas durante cada iteración.
    
    \item \textbf{Equipo de Desarrollo}
    \newline
    Este rol ha sido asignado al autor del proyecto debido a que ha sido el
    responsable de la ejecución técnica y el desarrollo del software.
\end{itemize}

\subsection{Artefactos de Scrum}

Los artefactos, dentro de la metodología Scrum, es información que el equipo utiliza para definir el trabajo que hay que
realizar para crear un producto. Hay tres tipos de artefactos: \textit{backlog} del producto, \textit{backlog} del sprint y
meta o incremento del sprint.

\begin{itemize}
    \item \textbf{Pila (\textit{backlog}) del Producto}
    \newline
    En este artefacto se agrupan todas las actividades en forma de lista que un equipo tiene que completar para crear el 
    producto. En este proyecto se ha utilizado la herramienta de GitHub Projects para crear un tablero donde se recojan 
    todas las actividades que se han planteado (historias de usuario, requisitos, mejoras y correcciones).

    \item \textbf{\textit{Backlog} del Sprint}
    \newline
    Esta sería la lista de elementos (historias de usuario y actividades) que el equipo de desarrollo ha seleccionado para
    realizar durante un sprint. Con la herramienta de GitHub Projects, además de poder tener una pila de producto con todas
    las actividades, se pueden agrupar las actividades por iteraciones / sprints y crear tableros que solo visualicen las 
    actividades relacionadas con ese sprint.

    \item \textbf{Meta del Sprint}
    \newline
    Este artefacto corresponde a la parte del producto que se obtiene al acabar una iteración. En este caso, cada vez 
    que se acaba la iteración, se obtiene una nueva funcionalidad que se añade al producto y que puede ser probada, con el fin
    de ir obteniendo algo ``tangible'' que el \textit{Product Owner} y el cliente pueda ver.
\end{itemize}

\subsection{Eventos de Scrum}

Con eventos de Scrum se va a referir a todas aquellas prácticas que son realizadas dentro del equipo Scrum de manera periódica,
englobando las ceremonias y reuniones que estos realizan en distintos momentos durante el desarrollo del producto.

Entre los eventos de Scrum, se encuentran los siguientes:

\begin{itemize}
    \item \textbf{Organización del Backlog}
    \newline
    Consiste en realizar una limpieza y actualización sobre la pila del producto, de manera que pueda estar lista para
    usarse por parte del equipo de desarrollo y de los miembros del equipo en cualquier momento. Esta labor es realizada
    por el \textit{Product Owner}.

    \item \textbf{Planificación de Sprints}
    \newline
    La planificación del sprint consiste en una reunión que realiza el equipo de desarrollo para planificar cual será
    el trabajo que se va a realizar durante el sprint. Esta reunión es diriga por el \textit{Scrum Master} y va a servir
    para ir añadiendo las historias de usuario o actividades importantes que se tengan que realizar, pasándolas del
    \textit{Product Backlog} al sprint específico.

    \item \textbf{Ejecución de Sprint}
    \newline
    Esta fase corresponde a la ejecución de todo lo planteado en el sprint, desarrollando e implementando todas aquellas
    tareas que han sido fijadas.

    \item \textbf{Reunión rápida}
    \newline
    Este tipo de ruenión se realiza al comienzo de la jornada de manera diaria. Son reuniones de corta duración donde se
    plantea lo que se va a realizar ese día y si hay algún obstáculo, de manera que todo el equipo trabaje en sintonía y
    se acerquen cada vez más a la consecución del sprint.

    \item \textbf{Revisión del Sprint}
    \newline
    La revisión del sprint consiste en una reunión en la que el equipo de desarrollo muestra a las partes interesadas del
    producto los elementos del backlog que están finalizados a través de una ``demo'', diciendo el \textit{Product Owner}
    si se publica o no el incremento.

    \item \textbf{Retrospectiva del Sprint}
    \newline
    Esta es la última reunión, donde el equipo se reúne para analizar qué cosas se han conseguido y han funcionado y cuales
    no, buscando de esta forma qué cosas deben mejorarse de cara a las próximas iteraciones.
\end{itemize}

\subsection{¿Por qué usar la metodología Scrum?}



\subsection{Adapatación de la metodología Scrum al proyecto}

Una vez explicadas todas las partes que componen la metodología Scrum, se va a proceder a explicar cómo esas partes han
sido adaptadas y tratadas en el presente proyecto.

Como se ha comentado anteriormente, se ha utilizado la herramienta de GitHub Projects \cite{githubProjects} para elaborar
el \textbf{\textit{Product Backlog}}, donde se han añadido todas las historias de usuario así como actividades necesarias
para poder desarrollar el software. A través de esta pila de producto se han ido organizando los diferentes sprints con 
las actividades necesarias para la realización de cada uno de ellos, ya que la herramienta de Projects permite crear distintas
vistas (cada una corresponde a un sprint) y asignar tareas del backlog a estas vistas. 

Tanto el Product Backlog como la planificación de cada uno de los sprints han sido llevados a cabo por el autor del proyecto, 
añadiendo toda la información (historias de usuario y tareas) que ha considerado necesaria para desarrollar el software 
junto con las recomendaciones del tutor del TFG (Dr. Juan Luis Jiménez Laredo) cuando ha sido requerido. Asimismo, la 
ejecución de cada sprint ha sido realizada por el autor del proyecto. En las imágenes ... se puede ver un ejemplo de cómo es
una historia de usurio y de cómo sería una sprint:

INSERTAR IMAGEN DE HISTORIA USUARIO

INSERTAR IMAGEN DE SPRINT 1

En cuanto a las reuniones realizadas durante todo el desarrollo del sistema, han sido tratas de la siguiente forma:

\begin{itemize}
    \item \textbf{Reunión diaria}: ha sido realizada por el autor del proyecto, revisando qué hizo el último día y planteando
    las siguientes tareas a desarrollar, teniendo en cuenta los plazos fijados y reorganizando tanto el sprint como
    el backlog en caso necesario. 

    \item \textbf{Revisión del sprint}: esta reunión ha sido realizada entre el tutor del TFG y el autor del proyecto al finalizar
    cada sprint. En esta reunión se han explicado las partes del software que han sido desarrolladas, se ha mostrado el funcionamiento
    del sistema y se ha debatido sobre posibles correcciones a realizar.

    \item \textbf{Reunión de retrospectiva}: ha sido realizada junto con el tutor del TFG y ha servido para tratar qué aspectos de lo
    realizado durante el sprint es mejorable y qué se podría añadir o contemplar para futuros sprints. También ha servido para fijar las 
    ideas de lo que iba a ser desarrollado en el próximo sprint y qué camino se iba a ir tomando.
\end{itemize}

\subsection{Sprints realizados}

Hablar de cada uno de los sprints realizados y del cronograma

\section{Presupuesto del proyecto}